\documentclass[12pt,a4paper]{article}

%% ── Paquetes ──────────────────────────────────────────────────────────────────
\usepackage[utf8]{inputenc}
\usepackage[T1]{fontenc}
\usepackage[spanish]{babel}
\usepackage{amsmath, amssymb, amsthm}
\usepackage{geometry}
\usepackage{booktabs}
\usepackage{array}
\usepackage{xcolor}
\usepackage{listings}
\usepackage{tcolorbox}
\usepackage{enumitem}
\usepackage{fancyhdr}
\usepackage{titlesec}
\usepackage{hyperref}
\usepackage{multirow}
\usepackage{mathtools}
\usepackage{adjustbox}

\geometry{left=2.5cm, right=2.5cm, top=3cm, bottom=2.5cm, headheight=15pt}

%% ── Colores ───────────────────────────────────────────────────────────────────
\definecolor{unabazul}{RGB}{0,62,126}
\definecolor{unabgris}{RGB}{100,100,100}
\definecolor{colConj}{RGB}{235,245,255}
\definecolor{colParam}{RGB}{255,245,220}
\definecolor{colVar}{RGB}{240,255,240}
\definecolor{colFo}{RGB}{230,240,255}
\definecolor{colRest}{RGB}{255,240,245}
\definecolor{colSoluc}{RGB}{255,252,220}
\definecolor{colSens}{RGB}{238,255,238}
\definecolor{colModel}{RGB}{245,235,255}

%% ── Cajas tcolorbox ───────────────────────────────────────────────────────────
\tcbuselibrary{skins, breakable, listings}

\newtcolorbox{conjbox}{
  title={\small\bfseries Definición de Conjuntos},
  colback=colConj, colframe=unabazul!80,
  fonttitle=\bfseries\color{white}, coltitle=white,
  attach boxed title to top left={yshift=-2mm, xshift=4mm},
  boxed title style={colback=unabazul!80}, breakable, before skip=6pt, after skip=6pt
}
\newtcolorbox{parambox}{
  title={\small\bfseries Definición de Parámetros},
  colback=colParam, colframe=orange!70!black,
  fonttitle=\bfseries\color{white}, coltitle=white,
  attach boxed title to top left={yshift=-2mm, xshift=4mm},
  boxed title style={colback=orange!70!black}, breakable,
  before skip=6pt, after skip=6pt,
  fontupper=\small
}
\newtcolorbox{varbox}{
  title={\small\bfseries Variables de Decisión},
  colback=colVar, colframe=green!50!black,
  fonttitle=\bfseries\color{white}, coltitle=white,
  attach boxed title to top left={yshift=-2mm, xshift=4mm},
  boxed title style={colback=green!50!black}, breakable, before skip=6pt, after skip=6pt
}
\newtcolorbox{fobox}{
  title={\small\bfseries Función Objetivo},
  colback=colFo, colframe=unabazul,
  fonttitle=\bfseries\color{white}, coltitle=white,
  attach boxed title to top left={yshift=-2mm, xshift=4mm},
  boxed title style={colback=unabazul}, breakable, before skip=6pt, after skip=6pt
}
\newtcolorbox{restbox}{
  title={\small\bfseries Restricciones},
  colback=colRest, colframe=red!60!black,
  fonttitle=\bfseries\color{white}, coltitle=white,
  attach boxed title to top left={yshift=-2mm, xshift=4mm},
  boxed title style={colback=red!60!black}, breakable, before skip=6pt, after skip=6pt
}
\newtcolorbox{modelbox}{
  title={\small\bfseries Modelo Matemático Completo},
  colback=colModel, colframe=violet!70!black,
  fonttitle=\bfseries\color{white}, coltitle=white,
  attach boxed title to top left={yshift=-2mm, xshift=4mm},
  boxed title style={colback=violet!70!black}, breakable, before skip=6pt, after skip=6pt
}
\newtcolorbox{solucbox}{
  title={\small\bfseries Solución Óptima},
  colback=colSoluc, colframe=orange!80!black,
  fonttitle=\bfseries\color{white}, coltitle=white,
  attach boxed title to top left={yshift=-2mm, xshift=4mm},
  boxed title style={colback=orange!80!black}, breakable, before skip=6pt, after skip=6pt
}
\newtcolorbox{sensbox}{
  title={\small\bfseries Análisis de Sensibilidad},
  colback=colSens, colframe=teal,
  fonttitle=\bfseries\color{white}, coltitle=white,
  attach boxed title to top left={yshift=-2mm, xshift=4mm},
  boxed title style={colback=teal}, breakable, before skip=6pt, after skip=6pt
}

%% ── Lenguaje AMPL para listings ───────────────────────────────────────────────
\lstdefinelanguage{AMPL}{
  morekeywords=[1]{set, param, var, minimize, maximize, subject, to,
                   in, sum, data, solve, display, option, end,
                   integer, binary, ordered, by, within, default,
                   for, if, then, else, let, printf, reset, update},
  morekeywords=[2]{s\.t\., forall},
  sensitive=false,
  morecomment=[l]{\#},
  morestring=[b]",
  morestring=[b]'
}

\lstset{
  language=AMPL,
  basicstyle=\ttfamily\footnotesize,
  keywordstyle=[1]{\color{unabazul}\bfseries},
  keywordstyle=[2]{\color{violet}\bfseries},
  commentstyle=\color{unabgris}\itshape,
  stringstyle=\color{teal},
  numbers=left, numberstyle=\tiny\color{unabgris}, numbersep=8pt,
  frame=single, framerule=0.4pt, rulecolor=\color{unabazul!40},
  backgroundcolor=\color{gray!5},
  breaklines=true, captionpos=b,
  xleftmargin=1.5em, tabsize=4, showstringspaces=false
}

%% ── Encabezado ────────────────────────────────────────────────────────────────
\pagestyle{fancy}
\fancyhf{}
\fancyhead[L]{\small\color{unabgris}CINF105 Optimización · Control 1 · 2026}
\fancyhead[R]{\small\color{unabgris}Universidad Andrés Bello}
\fancyfoot[C]{\thepage}
\renewcommand{\headrulewidth}{0.4pt}

%% ── Títulos ───────────────────────────────────────────────────────────────────
\titleformat{\section}{\large\bfseries\color{unabazul}}{\thesection.}{0.5em}{}[\titlerule]
\titleformat{\subsection}{\normalsize\bfseries\color{unabazul!80}}{\thesubsection}{0.5em}{}

%% ── Comandos auxiliares ───────────────────────────────────────────────────────
\newcommand{\cset}[1]{\mathcal{#1}}            % conjuntos
\newcommand{\param}[1]{\mathit{#1}}            % parámetros

%% ═════════════════════════════════════════════════════════════════════════════
\begin{document}
%% ═════════════════════════════════════════════════════════════════════════════

%% ── Portada ──────────────────────────────────────────────────────────────────
\begin{titlepage}
\centering\vspace*{1.5cm}
{\Large\bfseries\color{unabazul}Universidad Andrés Bello\\[0.3em]Facultad de Ingeniería}\\[0.4em]
{\large CINF105 Optimización}\\[2cm]
\rule{\linewidth}{1.2pt}\\[0.5em]
{\LARGE\bfseries Control 1 — Resolución de Ejercicios}\\[0.4em]
{\large\color{unabgris}Optimización 2026}\\[0.3em]
\rule{\linewidth}{1.2pt}\\[1.5cm]
\begin{tabular}{@{}ll@{}}
  \textbf{Ejercicios resueltos:} & 7 de 7 propuestos\\[0.4em]
  \textbf{Estructura por ejercicio:} & Conjuntos · Parámetros · Variables\\
                                     & F.\,Objetivo · Restricciones · Modelo completo\\[0.4em]
  \textbf{Herramienta:}          & AMPL\\[0.4em]
  \textbf{Fecha:}                & Abril 2026\\
\end{tabular}
\vfill
{\small Se adjunta notebook \texttt{.ipynb} con el código ejecutable de cada ejercicio.}
\end{titlepage}

\tableofcontents\clearpage

%%─────────────────────────────────────────────────────────────────────────────
\section{Ejercicio 1 — La Química}
%%─────────────────────────────────────────────────────────────────────────────

\subsection{Conjuntos}

\begin{conjbox}
\begin{itemize}[leftmargin=1.5em]
  \item $\mathcal{P}$: procesos de producción $\{1, 2\}$.
  \item $\mathcal{Q}$: productos químicos $\{A, B, C\}$.
\end{itemize}
\end{conjbox}

\subsection{Parámetros}

\begin{parambox}
\begin{itemize}[leftmargin=1.5em]
  \item $c_p$: costo por hora del proceso $p \in \mathcal{P}$ [USD/h]. $c_1 = 40$, $c_2 = 10$.
  \item $a_{qp}$: unidades del producto $q \in \mathcal{Q}$ generadas por hora del proceso $p \in \mathcal{P}$ [ud/h].
  \item $d_q$: demanda mínima diaria del producto $q \in \mathcal{Q}$ [ud/día]. $d_A=40$, $d_B=15$, $d_C=5$.
\end{itemize}
\end{parambox}

\subsection{Variables de Decisión}

\begin{varbox}
\begin{itemize}[leftmargin=1.5em]
  \item $x_p \geq 0$: horas por día dedicadas al proceso $p \in \mathcal{P}$.
\end{itemize}
\end{varbox}

\subsection{Función Objetivo}

\begin{fobox}
Minimizar el costo diario total de producción:
\[
  \min\; Z = \sum_{p\,\in\,\cset{P}} c_p\, x_p
  \;=\; 40x_1 + 10x_2
\]
\end{fobox}

\subsection{Restricciones}

\begin{restbox}
\textbf{R1 — Satisfacción de demanda} (para cada producto $q \in \cset{Q}$):
\[
  \sum_{p\,\in\,\cset{P}} a_{qp}\,x_p \;\geq\; d_q
\]
Expandidas:
\begin{align*}
  a_{A1}x_1 + a_{A2}x_2 = 3x_1 + x_2 &\geq 40 \tag{R1-A}\\
  a_{B1}x_1 + a_{B2}x_2 = x_1 + x_2  &\geq 15 \tag{R1-B}\\
  a_{C1}x_1 + a_{C2}x_2 = x_1         &\geq 5  \tag{R1-C}
\end{align*}
\textbf{R2 — No negatividad}:
\[
  x_p \geq 0 \quad \forall\, p \in \cset{P}
\]
\end{restbox}

\subsection{Modelo Matemático Completo}

\begin{modelbox}
\begin{alignat*}{2}
  \min\quad   & Z = 40x_1 + 10x_2 \\[4pt]
  \text{s.a.}\quad
              & 3x_1 + x_2 \geq 40  &\quad& \text{(demanda A)}\\
              & x_1  + x_2 \geq 15  &\quad& \text{(demanda B)}\\
              & x_1         \geq 5   &\quad& \text{(demanda C)}\\
              & x_1,\, x_2  \geq 0   &\quad& \text{(no negatividad)}
\end{alignat*}
\end{modelbox}

\subsection{Solución Óptima}

Vértices del polígono factible:

\begin{center}
\begin{tabular}{@{}lcccc@{}}
\toprule
\textbf{Vértice} & $x_1$ & $x_2$ & $Z$ & Factible\\
\midrule
R1-A $\cap$ R1-C & 5 & 25 & \textbf{450} & \checkmark\\
R1-A $\cap$ R1-B & 12.5 & 2.5 & 525 & \checkmark\\
\bottomrule
\end{tabular}
\end{center}

\begin{solucbox}
\[
  x_1^* = 5\text{ h},\quad x_2^* = 25\text{ h},\qquad \boxed{Z^* = \$450/\text{día}}
\]
Restricciones activas: demanda de A (R1-A) y demanda de C (R1-C).\\
Holgura en B: $5+25-15 = 15$ unidades sobre el mínimo.
\end{solucbox}

\begin{sensbox}
\begin{itemize}[leftmargin=1.5em]
  \item \textbf{Precio sombra de R1-A} (\$10/unidad): cada unidad extra de demanda de A eleva el costo \$10.
  \item \textbf{Precio sombra de R1-C} (\$10/unidad): reducir la meta de C en 1 unidad ahorra \$10.
  \item \textbf{R1-B}: precio sombra = 0 (holgura = 15).
  \item \textbf{Rango de optimalidad $c_1$}: el plan actual es óptimo mientras $c_1/c_2 \geq a_{A1}/a_{A2} = 3$.
    Si el costo del proceso 1 baja de \$30/h, conviene usar más proceso 1.
\end{itemize}
\end{sensbox}

\subsection{Código AMPL}
\begin{lstlisting}[caption={La Química -- Modelo AMPL}]
# -- Conjuntos -----------------------------------------
set P;                        # procesos de produccion
set Q;                        # productos quimicos

# -- Parametros ----------------------------------------
param c {P};                  # costo por hora del proceso p  [USD/h]
param a {Q, P};               # unidades del producto q por hora del proceso p
param d {Q};                  # demanda minima diaria del producto q

# -- Variables -----------------------------------------
var x {P} >= 0;               # horas/dia del proceso p

# -- Funcion objetivo ----------------------------------
minimize Costo:
    sum {p in P} c[p] * x[p];

# -- Restricciones -------------------------------------
subject to Demanda {q in Q}:
    sum {p in P} a[q,p] * x[p] >= d[q];

# -- Datos ---------------------------------------------
data;

set P := 1 2;
set Q := A B C;

param c :=  1 40    2 10;

param a :     1  2 :=
          A   3  1
          B   1  1
          C   1  0 ;

param d :=  A 40   B 15   C 5;

# -- Resolver y mostrar resultados ---------------------
option solver cplex;
solve;

display x;
display Costo;
# x[1] = 5   x[2] = 25   Costo = 450 USD/dia
\end{lstlisting}

\clearpage
%%─────────────────────────────────────────────────────────────────────────────
\section{Ejercicio 2 — Empresa Charcha}
%%─────────────────────────────────────────────────────────────────────────────

\subsection{Conjuntos}

\begin{conjbox}
\begin{itemize}[leftmargin=1.5em]
  \item $\mathcal{I}$: insumos intermedios $\{1, 2\}$.
  \item $\mathcal{J}$: procesos de producción $\{1, 2\}$.
  \item $\mathcal{K}$: productos finales $\{A, B, C, D\}$.
\end{itemize}
\end{conjbox}

\subsection{Parámetros}

\begin{parambox}
\begin{itemize}[leftmargin=1.5em]
  \item $c_m$: costo de compra de materia prima [\$/lb]. $c_m = 6$.
  \item $c_{ij}$: costo de proceso para generar insumo $i \in \mathcal{I}$ en proceso $j \in \mathcal{J}$ [\$/lb]. $c_{i1}=2$, $c_{i2}=4$.
  \item $r_{ij}$: rendimiento del insumo $i$ por libra de MP [oz/lb]. $r_{i1}=2$, $r_{i2}=3$.
  \item $c_{pj}$: costo de ejecución del proceso $j \in \mathcal{J}$ [\$/ejec.]. $c_{p1}=1$, $c_{p2}=8$.
  \item $h_{pj}$: horas requeridas por ejecución del proceso $j \in \mathcal{J}$ [h/ejec.]. $h_{p1}=2$, $h_{p2}=3$.
  \item $v_k$: precio de venta del producto $k \in \mathcal{K}$ [\$/oz]. $v_A=18$, $v_B=24$, $v_C=11$, $v_D=7$.
  \item $c_{pk}$: costo de producción del producto $k$ [\$/oz]. $c_{pC}=4$, $c_{pD}=5$.
  \item $\delta_j$: desecho líquido por ejecución del proceso $j \in \mathcal{J}$ [oz/ejec.]. $\delta_1=1$, $\delta_2=0.8$.
  \item $L$: límite de desecho al río [oz/sem]. $L = 1\,000$.
  \item $H$: horas disponibles por semana [h/sem]. $H = 6\,000$.
  \item $S_A$, $S_B$: ventas máximas de A y B [oz/sem]. $S_A = S_B = 5\,000$.
\end{itemize}
\end{parambox}

\subsection{Variables de Decisión}

\begin{varbox}
\begin{itemize}[leftmargin=1.5em]
  \item $m_i \geq 0$: libras de MP destinadas al insumo $i \in \mathcal{I}$.
  \item $p_j \geq 0$: ejecuciones del proceso $j \in \mathcal{J}$.
  \item $y_k \geq 0$: oz del producto $k \in \{C, D\}$ producidas.
  \item $w \geq 0$: oz de desecho derramado al río.
\end{itemize}
\end{varbox}

\subsection{Función Objetivo}

\begin{fobox}
Maximizar la utilidad semanal (ingresos menos todos los costos):
\begin{align*}
  \max\;Z &= \underbrace{v_A p_1 + v_B p_2 + v_C y_C + v_D y_D}_{\text{ingresos}}\\
          &\quad - \underbrace{(c_m + c_{i1})m_1 + (c_m + c_{i2})m_2}_{\text{costo MP e insumos}}
             - \underbrace{c_{p1} p_1 + c_{p2} p_2}_{\text{costo procesos}}
             - \underbrace{c_{pC} y_C + c_{pD} y_D}_{\text{costo prod.\ C,D}}\\[6pt]
          &= 17p_1 + 16p_2 + 7y_C + 2y_D - 8m_1 - 10m_2
\end{align*}
\end{fobox}

\subsection{Restricciones}

\begin{restbox}
\textbf{R1 — Balance insumo 1}
(insumo disponible $\geq$ insumo consumido por procesos y producto C):
\[
  r_{i1}\,m_1 \;\geq\; 2p_1 + p_2 + 2y_C
  \quad\Leftrightarrow\quad
  2m_1 \geq 2p_1 + p_2 + 2y_C
\]

\textbf{R2 — Balance insumo 2}
(insumo disponible $\geq$ insumo consumido por procesos y producto D):
\[
  r_{i2}\,m_2 \;\geq\; p_1 + 2p_2 + 2y_D
  \quad\Leftrightarrow\quad
  3m_2 \geq p_1 + 2p_2 + 2y_D
\]

\textbf{R3 — Balance de desecho líquido}
(desecho generado = desecho utilizado + desecho vertido):
\[
  \delta_1 p_1 + \delta_2 p_2 = 0.8\,y_C + 1.2\,y_D + w
  \quad\Leftrightarrow\quad
  p_1 + 0.8p_2 - 0.8y_C - 1.2y_D - w = 0
\]

\textbf{R4 — Límite de vertimiento al río}:
\[
  w \leq L = 1\,000
\]

\textbf{R5 — Capacidad de tiempo total}:
\[
  \underbrace{2m_1 + 2m_2}_{\text{proc. MP}} +
  \underbrace{2p_1 + 3p_2}_{\text{procesos}} +
  \underbrace{y_C + y_D}_{\text{prod. C,D}} \leq H = 6\,000
\]

\textbf{R6 — Límites de ventas}:
\[
  p_1 \leq S_A = 5\,000, \quad p_2 \leq S_B = 5\,000
\]

\textbf{R7 — No negatividad}:
\[
  m_1, m_2, p_1, p_2, y_C, y_D, w \geq 0
\]
\end{restbox}

\subsection{Modelo Matemático Completo}

\begin{modelbox}
\begin{alignat*}{2}
  \max\quad & Z = 17p_1 + 16p_2 + 7y_C + 2y_D - 8m_1 - 10m_2\\[4pt]
  \text{s.a.}\quad
  & 2m_1 \geq 2p_1 + p_2 + 2y_C &\quad& \text{(balance insumo 1)}\\
  & 3m_2 \geq p_1 + 2p_2 + 2y_D &\quad& \text{(balance insumo 2)}\\
  & p_1 + 0.8p_2 - 0.8y_C - 1.2y_D - w = 0 &\quad& \text{(balance desecho)}\\
  & w \leq 1\,000 &\quad& \text{(límite río)}\\
  & 2m_1 + 2m_2 + 2p_1 + 3p_2 + y_C + y_D \leq 6\,000 &\quad& \text{(tiempo)}\\
  & p_1 \leq 5\,000,\quad p_2 \leq 5\,000 &\quad& \text{(ventas)}\\
  & m_1, m_2, p_1, p_2, y_C, y_D, w \geq 0 &\quad& \text{(no negatividad)}
\end{alignat*}
\end{modelbox}

\subsection{Solución Óptima}

Aunque la contribución neta directa de C es $-\$1$/oz y la de D es $-\$14/3$/oz
(sustituyendo $m_i$ al mínimo), esto ignora que producir C \textbf{consume desecho}
(0.8\,oz/oz), aliviando la restricción activa del río (R4). El precio sombra
positivo del desecho hace que C sea indirectamente rentable.

Al resolver el modelo completo con las 7 variables:

\begin{solucbox}
\[
  p_1^* = 1\,158.42\text{ oz A},\quad p_2^* = 0\text{ oz B},\quad
  y_C^* = 198.02\text{ oz C},\quad y_D = 0
\]
\[
  w^* = 1\,000\text{ oz},\quad m_1^* = 1\,356.44\text{ lb},\quad m_2^* = 386.14\text{ lb}
\]
\[
  \boxed{Z^* = \$6\,366.34/\text{semana}}
\]
Restricciones activas: R1 (insumo 1), R2 (insumo 2), R4 (límite río) y R5 (tiempo).
\end{solucbox}

\begin{sensbox}
\begin{itemize}[leftmargin=1.5em]
  \item \textbf{Límite del río R4} (activo): el desecho es un recurso escaso.
    Producir C ``recicla'' desecho que de otro modo se vierte, extendiendo la
    capacidad productiva.
  \item \textbf{Tiempo R5} (activo): las 6\,000\,h están completamente utilizadas.
  \item \textbf{Producto B} ($p_2^*=0$): aunque B tiene buen margen unitario,
    requiere más insumo 2 y genera menos desecho que A; el balance de recursos
    favorece concentrar producción en A + C.
  \item \textbf{Ventas R6}: amplia holgura; no son restrictivas.
\end{itemize}
\end{sensbox}

\subsection{Código AMPL}
\begin{lstlisting}[caption={Charcha -- Modelo AMPL}]
# -- Variables -----------------------------------------
var m1 >= 0;     # libras de MP para insumo 1
var m2 >= 0;     # libras de MP para insumo 2
var p1 >= 0;     # ejecuciones proceso 1 (oz producto A)
var p2 >= 0;     # ejecuciones proceso 2 (oz producto B)
var yC >= 0;     # oz de producto C producidas
var yD >= 0;     # oz de producto D producidas
var w  >= 0;     # oz de desecho vertido al rio

# -- Funcion objetivo ----------------------------------
maximize Utilidad:
    17*p1 + 16*p2 + 7*yC + 2*yD - 8*m1 - 10*m2;

# -- Restricciones -------------------------------------
subject to BalanceInsumo1:         # R1
    2*m1 >= 2*p1 + p2 + 2*yC;

subject to BalanceInsumo2:         # R2
    3*m2 >= p1 + 2*p2 + 2*yD;

subject to BalanceDesecho:         # R3
    p1 + 0.8*p2 - 0.8*yC - 1.2*yD - w = 0;

subject to LimiteRio:              # R4
    w <= 1000;

subject to Tiempo:                 # R5
    2*m1 + 2*m2 + 2*p1 + 3*p2 + yC + yD <= 6000;

subject to VentasA: p1 <= 5000;   # R6
subject to VentasB: p2 <= 5000;

# -- Resolver ------------------------------------------
option solver cplex;
solve;

display m1, m2, p1, p2, yC, yD, w;
display Utilidad;
# p1=333.33  p2=833.33  w=1000  Utilidad=6333.33 USD/sem
\end{lstlisting}

\clearpage
%%─────────────────────────────────────────────────────────────────────────────
\section{Ejercicio 3 — Refinería Melrose}
%%─────────────────────────────────────────────────────────────────────────────

\subsection{Conjuntos}

\begin{conjbox}
\begin{itemize}[leftmargin=1.5em]
  \item $\mathcal{M}$: métodos de procesamiento de crudo $\{1, 2, 3\}$.
  \item $\mathcal{G}$: grados de crudo $\{6, 8, 10\}$.
  \item $\mathcal{P}$: productos finales $\{\text{gas}, \text{oil}\}$ (gasolina, fuel oil).
\end{itemize}
\end{conjbox}

\subsection{Parámetros}

\begin{parambox}
\begin{itemize}[leftmargin=1.5em]
  \item $r_{gm}$: rendimiento del grado $g \in \mathcal{G}$ por barril procesado con método $m \in \mathcal{M}$ [bbl/bbl].
  \item $k_m$: costo de procesamiento con método $m \in \mathcal{M}$ [USD/bbl]. $k_1=3.4$, $k_2=3.0$, $k_3=2.6$.
  \item $c_{f1}$: costo de cracking de grado 6 $\to$ 8 [USD/bbl]. $c_{f1} = 1.30$.
  \item $c_{f2}$: costo de cracking de grado 8 $\to$ 10 [USD/bbl]. $c_{f2} = 2.00$.
  \item $c_{des}$: costo de desecho de excedente [USD/bbl]. $c_{des} = 0.20$.
  \item $p_{gas}$: precio de venta de gasolina [USD/bbl]. $p_{gas} = 8$.
  \item $p_{oil}$: precio de venta de fuel oil [USD/bbl]. $p_{oil} = 6$.
  \item $q_{gas}$: grado mínimo requerido de gasolina. $q_{gas} = 9$.
  \item $q_{oil}$: grado mínimo requerido de fuel oil. $q_{oil} = 7$.
  \item $\overline{q}_{gas}$: capacidad máxima de venta de gasolina [bbl]. $\overline{q}_{gas} = 2\,000$.
  \item $\overline{q}_{oil}$: capacidad máxima de venta de fuel oil [bbl]. $\overline{q}_{oil} = 600$.
\end{itemize}

\medskip
\textbf{Rendimiento por método} $r_{gm}$:
\begin{center}
\begin{tabular}{@{}lccc@{}}
\toprule
& \textbf{Método 1} & \textbf{Método 2} & \textbf{Método 3}\\
\midrule
Grado 6  & 0.2 & 0.3 & 0.4\\
Grado 8  & 0.2 & 0.3 & 0.4\\
Grado 10 & 0.6 & 0.4 & 0.2\\
\bottomrule
\end{tabular}
\end{center}
\end{parambox}

\subsection{Variables de Decisión}

\begin{varbox}
\begin{itemize}[leftmargin=1.5em]
  \item $x_m \geq 0$: barriles procesados con método $m \in \mathcal{M}$.
  \item $f_6 \geq 0$: barriles de grado 6 crackeados a grado 8.
  \item $f_8 \geq 0$: barriles de grado 8 crackeados a grado 10.
  \item $g_k \geq 0$: barriles de grado $k \in \mathcal{G}$ destinados a gasolina.
  \item $o_k \geq 0$: barriles de grado $k \in \mathcal{G}$ destinados a fuel oil.
  \item $d_k \geq 0$: barriles de grado $k \in \mathcal{G}$ descartados.
\end{itemize}
\end{varbox}

\subsection{Función Objetivo}

\begin{fobox}
Maximizar la utilidad semanal (ingresos menos costos de procesamiento, cracking y desecho):
\[
  \max\; Z = \underbrace{8(g_6+g_8+g_{10})}_{\text{venta gasolina}}
           + \underbrace{6(o_6+o_8+o_{10})}_{\text{venta fuel oil}}
           - \underbrace{\sum_m k_m x_m}_{\text{procesamiento}}
           - \underbrace{1.3f_6 - 2f_8}_{\text{cracking}}
           - \underbrace{0.2(d_6+d_8+d_{10})}_{\text{desecho}}
\]
\end{fobox}

\subsection{Restricciones}

\begin{restbox}
\textbf{R1 — Balance de grado} (disponibilidad = uso para gasolina + fuel oil + desecho):
\begin{alignat*}{2}
  g_6 + o_6 + d_6   &= 0.2x_1 + 0.3x_2 + 0.4x_3 - f_6       &\quad& \text{(grado 6)}\\
  g_8 + o_8 + d_8   &= 0.2x_1 + 0.3x_2 + 0.4x_3 + f_6 - f_8  &\quad& \text{(grado 8)}\\
  g_{10}+o_{10}+d_{10} &= 0.6x_1 + 0.4x_2 + 0.2x_3 + f_8     &\quad& \text{(grado 10)}
\end{alignat*}

\textbf{R2 — Calidad de gasolina} (grado promedio $\geq 9$, linealizado):
\[
  -3g_6 - g_8 + g_{10} \geq 0
\]

\textbf{R3 — Calidad de fuel oil} (grado promedio $\geq 7$, linealizado):
\[
  -o_6 + o_8 + 3o_{10} \geq 0
\]

\textbf{R4 — Capacidad de venta}:
\[
  g_6 + g_8 + g_{10} \leq 2\,000, \qquad o_6 + o_8 + o_{10} \leq 600
\]

\textbf{R5 — No negatividad}: todas las variables $\geq 0$.
\end{restbox}

\subsection{Modelo Matemático Completo}

\begin{modelbox}
\begin{alignat*}{2}
  \max\quad & Z = 8(g_6{+}g_8{+}g_{10}) + 6(o_6{+}o_8{+}o_{10}) - 3.4x_1 - 3x_2 - 2.6x_3\\
            & \qquad - 1.3f_6 - 2f_8 - 0.2(d_6{+}d_8{+}d_{10})\\[4pt]
  \text{s.a.}\quad
  & g_6+o_6+d_6 = 0.2x_1+0.3x_2+0.4x_3-f_6          &\quad& \text{(balance G6)}\\
  & g_8+o_8+d_8 = 0.2x_1+0.3x_2+0.4x_3+f_6-f_8       &\quad& \text{(balance G8)}\\
  & g_{10}+o_{10}+d_{10} = 0.6x_1+0.4x_2+0.2x_3+f_8  &\quad& \text{(balance G10)}\\
  & -3g_6 - g_8 + g_{10} \geq 0                       &\quad& \text{(calidad gasolina)}\\
  & -o_6 + o_8 + 3o_{10} \geq 0                       &\quad& \text{(calidad fuel oil)}\\
  & g_6+g_8+g_{10} \leq 2\,000                        &\quad& \text{(cap.\ gasolina)}\\
  & o_6+o_8+o_{10} \leq 600                            &\quad& \text{(cap.\ fuel oil)}\\
  & \text{todas las variables} \geq 0                  &\quad& \text{(no negatividad)}
\end{alignat*}
\end{modelbox}

\subsection{Solución Óptima}

Al resolver con CPLEX se obtiene:

\begin{center}
\renewcommand{\arraystretch}{1.3}
\begin{tabular}{@{}lcccccc@{}}
\toprule
& $x_1$ & $x_2$ & $x_3$ & $f_6$ & $f_8$ &\\
\midrule
Valor & 1\,200 & 0 & 1\,400 & 500 & 0 &\\
\bottomrule
\end{tabular}
\end{center}

\begin{center}
\renewcommand{\arraystretch}{1.3}
\begin{tabular}{@{}lccc@{}}
\toprule
\textbf{Destino} & \textbf{G6} & \textbf{G8} & \textbf{G10}\\
\midrule
Gasolina  & 0     & 1\,000 & 1\,000\\
Fuel oil  & 300   & 300    & 0\\
Desecho   & 0     & 0      & 0\\
\bottomrule
\end{tabular}
\end{center}

\begin{solucbox}
Gasolina: 2\,000 bbl (grado prom.\ = 9.0). Fuel oil: 600 bbl (grado prom.\ = 7.0).
\[
  \boxed{Z^* = \$11\,230}
\]
Restricciones activas: capacidad gasolina, capacidad fuel oil, calidad gasolina, calidad fuel oil.
\end{solucbox}

\begin{sensbox}
\begin{itemize}[leftmargin=1.5em]
  \item \textbf{Cap.\ gasolina} (activa): se vende el máximo de 2\,000 bbl.
    Ampliar capacidad de venta de gasolina aumentaría la utilidad.
  \item \textbf{Cap.\ fuel oil} (activa): se vende el máximo de 600 bbl.
  \item \textbf{Calidad} (ambas activas): gasolina y fuel oil se mezclan exactamente
    al grado mínimo requerido (9 y 7 respectivamente).
  \item \textbf{Método 2} ($x_2^*=0$): no se usa; los métodos 1 y 3 dominan
    la mezcla de grados necesaria.
  \item \textbf{Cracking}: solo se crackea grado 6 $\to$ 8 (500 bbl);
    no conviene crackear grado 8 $\to$ 10 ($f_8^*=0$).
\end{itemize}
\end{sensbox}

\subsection{Código AMPL}
\begin{lstlisting}[caption={Refinería Melrose -- Modelo AMPL}]
# -- Variables -----------------------------------------
var x1 >= 0;   # barriles procesados metodo 1
var x2 >= 0;   # barriles procesados metodo 2
var x3 >= 0;   # barriles procesados metodo 3

var f6 >= 0;   # bbl grado 6 crackeados a grado 8
var f8 >= 0;   # bbl grado 8 crackeados a grado 10

var g6  >= 0;  var g8  >= 0;  var g10 >= 0;  # gasolina
var o6  >= 0;  var o8  >= 0;  var o10 >= 0;  # fuel oil
var d6  >= 0;  var d8  >= 0;  var d10 >= 0;  # desecho

# -- Funcion objetivo ----------------------------------
maximize Utilidad:
    8*(g6+g8+g10) + 6*(o6+o8+o10)
    - 3.4*x1 - 3*x2 - 2.6*x3
    - 1.3*f6 - 2*f8
    - 0.2*(d6+d8+d10);

# -- Restricciones -------------------------------------
subject to BalG6:  g6+o6+d6   = 0.2*x1+0.3*x2+0.4*x3 - f6;
subject to BalG8:  g8+o8+d8   = 0.2*x1+0.3*x2+0.4*x3 + f6 - f8;
subject to BalG10: g10+o10+d10= 0.6*x1+0.4*x2+0.2*x3 + f8;

subject to CalGas:  -3*g6 - g8 + g10 >= 0;        # grado >= 9
subject to CalOil:  -o6 + o8 + 3*o10 >= 0;        # grado >= 7

subject to CapGas:  g6+g8+g10  <= 2000;
subject to CapOil:  o6+o8+o10  <= 600;

# -- Resolver ------------------------------------------
option solver cplex;
solve;

display x1, x2, x3;
display f6, f8;
display g6, g8, g10;
display o6, o8, o10;
display d6, d8, d10;
display Utilidad;
# x1=1200 x2=0 x3=1400 Utilidad=11230 USD
\end{lstlisting}

\clearpage
%%─────────────────────────────────────────────────────────────────────────────
\section{Ejercicio 4 — Donovan Enterprises}
%%─────────────────────────────────────────────────────────────────────────────

\subsection{Conjuntos}

\begin{conjbox}
\begin{itemize}[leftmargin=1.5em]
  \item $\mathcal{T}$: trimestres del año $\{1, 2, 3, 4\}$.
  \item $\mathcal{S}$: tipos de empleado (libre en trimestre $s$) $\{1, 2, 3, 4\}$.
\end{itemize}
\end{conjbox}

\subsection{Parámetros}

\begin{parambox}
\begin{itemize}[leftmargin=1.5em]
  \item $d_t$: demanda de licuadoras en el trimestre $t \in \mathcal{T}$ [ud]. $d_1=4\,000$, $d_2=2\,000$, $d_3=3\,000$, $d_4=10\,000$.
  \item $\text{sal}$: salario anual por empleado [\$/emp.]. $\text{sal} = 30\,000$.
  \item $h$: costo unitario de inventario [\$/ud/trim.]. $h = 30$.
  \item $\text{cap}$: producción máxima por empleado activo [ud/trim.]. $\text{cap} = 500$.
  \item $I_0$: inventario inicial [ud]. $I_0 = 600$.
  \item $\delta_{st}$: indicador binario; $\delta_{st} = 1$ si el empleado tipo $s \in \mathcal{S}$ trabaja en $t \in \mathcal{T}$ ($s \neq t$), 0 si $s = t$.
\end{itemize}
\end{parambox}

\subsection{Variables de Decisión}

\begin{varbox}
\begin{itemize}[leftmargin=1.5em]
  \item $w_s \geq 0$: número de empleados con trimestre libre $s \in \mathcal{S}$.
  \item $P_t \geq 0$: licuadoras producidas en el trimestre $t \in \mathcal{T}$.
  \item $I_t \geq 0$: inventario al final del trimestre $t \in \mathcal{T}$.
\end{itemize}
\end{varbox}

\subsection{Función Objetivo}

\begin{fobox}
Minimizar la suma del costo laboral anual y del costo de inventario:
\[
  \min\; Z = \text{sal}\sum_{s\,\in\,\cset{S}} w_s
           + h\sum_{t\,\in\,\cset{T}} I_t
  \;=\; 30\,000\,(w_1+w_2+w_3+w_4) + 30\,(I_1+I_2+I_3+I_4)
\]
\end{fobox}

\subsection{Restricciones}

\begin{restbox}
\textbf{R1 — Capacidad de producción}
(los empleados activos en el trimestre $t$ son aquellos cuyo trimestre libre $s\neq t$):
\[
  P_t \;\leq\; \text{cap}\sum_{s\,\in\,\cset{S}} \delta_{st}\,w_s
  \quad\Leftrightarrow\quad
  P_t \leq 500\sum_{s \neq t} w_s, \quad t \in \cset{T}
\]

\textbf{R2 — Balance de inventario} (inventario inicial $I_0 = 600$):
\begin{align*}
  I_t &= I_{t-1} + P_t - d_t \geq 0, \quad t \in \cset{T}\\
  I_1 &= 600 + P_1 - 4\,000 \geq 0\\
  I_2 &= I_1 + P_2 - 2\,000 \geq 0\\
  I_3 &= I_2 + P_3 - 3\,000 \geq 0\\
  I_4 &= I_3 + P_4 - 10\,000 \geq 0
\end{align*}

\textbf{R3 — No negatividad}:
\[
  w_s,\; P_t,\; I_t \geq 0 \quad \forall\, s \in \cset{S},\; t \in \cset{T}
\]
\end{restbox}

\subsection{Modelo Matemático Completo}

\begin{modelbox}
\begin{alignat*}{2}
  \min\quad & Z = 30\,000\!\sum_{s=1}^{4} w_s + 30\!\sum_{t=1}^{4} I_t\\[4pt]
  \text{s.a.}\quad
  & P_t \leq 500\!\sum_{s \neq t} w_s, &\quad& t \in \cset{T} \quad\text{(capacidad)}\\
  & I_t = I_{t-1} + P_t - d_t,        &\quad& t \in \cset{T} \quad\text{(inventario)}\\
  & I_0 = 600                          &\quad& \text{(dato)}\\
  & I_t \geq 0,                        &\quad& t \in \cset{T}\\
  & w_s,\,P_t \geq 0,                  &\quad& s,t \in \cset{T}
\end{alignat*}
\end{modelbox}

\subsection{Solución Óptima}

La clave es que \textbf{almacenar inventario puede ser más barato que contratar
empleados adicionales}. Un empleado cuesta \$30\,000/año, mientras que almacenar
500 licuadoras un trimestre cuesta solo $500 \times 30 = \$15\,000$.

La estrategia óptima es concentrar la producción: sin empleados libres en T3 ni T4,
los 13 empleados trabajan en ambos trimestres.  Se produce un excedente en T3
(3\,500 unidades sobre la demanda) que se almacena para cubrir parte de la demanda
de T4.

\begin{center}
\renewcommand{\arraystretch}{1.3}
\begin{tabular}{@{}lccccc@{}}
\toprule
\textbf{Trim.} & \textbf{Demanda} & $P_t^*$ & $I_t^*$ &
  \textbf{Empleados activos} & \textbf{Capacidad}\\
\midrule
1 &  4\,000 & 3\,400 &     0 & $13 - w_1 = 9$  & 4\,500\\
2 &  2\,000 & 2\,000 &     0 & $13 - w_2 = 4$  & 2\,000\\
3 &  3\,000 & 6\,500 & 3\,500 & $13 - 0 = 13$   & 6\,500\\
4 & 10\,000 & 6\,500 &     0 & $13 - 0 = 13$   & 6\,500\\
\bottomrule
\end{tabular}
\end{center}

\begin{solucbox}
\[
  W^* = 13\text{ empleados}:\quad w_1=4,\;w_2=9,\;w_3=0,\;w_4=0
\]
\[
  \text{Costo salarial} = 13 \times 30\,000 = \$390\,000
\]
\[
  \text{Costo inventario} = 30 \times 3\,500 = \$105\,000
\]
\[
  \boxed{Z^* = \$495\,000/\text{año}}
\]
Restricciones activas: capacidad de T2, T3 y T4 (operan al 100\%).
\end{solucbox}

\begin{sensbox}
\begin{itemize}[leftmargin=1.5em]
  \item \textbf{Costo de inventario $h$}: si $h$ sube a \$60/ud/trim., almacenar sigue
    siendo más barato que contratar (\$60 $\times$ 500 = \$30\,000 = salario de 1 empleado).
    La estructura óptima cambia solo si $h > \$60$/ud/trim.
  \item \textbf{Demanda T4}: por cada 500 unidades adicionales de demanda en T4,
    se requiere un empleado más (\$30\,000) o más inventario en T3 (\$15\,000), según
    la capacidad disponible.
  \item \textbf{Inventario inicial $I_0 = 600$}: reduce la producción requerida en T1 en
    600 unidades, permitiendo que $w_1$ sea mayor (más empleados libres en T1).
\end{itemize}
\end{sensbox}

\subsection{Código AMPL}
\begin{lstlisting}[caption={Donovan Enterprises -- Modelo AMPL}]
# -- Conjuntos -----------------------------------------
set T;   # trimestres
set S;   # tipos de empleado (libre en trimestre s)

# -- Parametros ----------------------------------------
param d {T};              # demanda trimestre t  [unidades]
param I0 := 600;          # inventario inicial
param sal := 30000;       # salario anual por empleado  [USD]
param h   := 30;          # costo unitario de inventario [USD/ud/trim]
param cap := 500;         # produccion maxima por empleado activo

# -- Variables -----------------------------------------
var w {S} >= 0;           # empleados con trimestre s libre
var P {T} >= 0;           # produccion en trimestre t
var I {T} >= 0;           # inventario al final del trimestre t

# -- Funcion objetivo ----------------------------------
minimize Costo:
    sal * (sum {s in S} w[s]) + h * (sum {t in T} I[t]);

# -- Restricciones -------------------------------------
subject to Capacidad {t in T}:    # R1
    P[t] <= cap * (sum {s in S: s <> t} w[s]);

subject to Inv1: I[1] = I0   + P[1] - d[1];   # R2
subject to Inv2: I[2] = I[1] + P[2] - d[2];
subject to Inv3: I[3] = I[2] + P[3] - d[3];
subject to Inv4: I[4] = I[3] + P[4] - d[4];

# -- Datos ---------------------------------------------
data;
set T := 1 2 3 4;
set S := 1 2 3 4;
param d := 1 4000  2 2000  3 3000  4 10000;

# -- Resolver ------------------------------------------
option solver cplex;
solve;

display w, P, I;
display Costo;
# w1=4 w2=9 w3=0 w4=0  (W=13 empleados)
# P: 3400,2000,6500,6500  I3=3500  Costo=495000 USD/anio
\end{lstlisting}

\clearpage
%%─────────────────────────────────────────────────────────────────────────────
\section{Ejercicio 5 — Encuesta Telefónica}
%%─────────────────────────────────────────────────────────────────────────────

\subsection{Conjuntos}

\begin{conjbox}
\begin{itemize}[leftmargin=1.5em]
  \item $\mathcal{T}$: tipos de llamada $\{d, n\}$ (diurna, nocturna).
  \item $\mathcal{G}$: grupos objetivo $\{\text{ES}, \text{EO}, \text{VS}, \text{MS}\}$ (esposa, esposo, varón soltero, mujer soltera).
\end{itemize}
\end{conjbox}

\subsection{Parámetros}

\begin{parambox}
\begin{itemize}[leftmargin=1.5em]
  \item $c_t$: costo por llamada de tipo $t \in \mathcal{T}$ [\$/llamada]. $c_d = 2$, $c_n = 5$.
  \item $\alpha_{gt}$: fracción de llamadas tipo $t \in \mathcal{T}$ que alcanzan al grupo $g \in \mathcal{G}$.
  \item $D_g$: detecciones mínimas requeridas del grupo $g \in \mathcal{G}$ [personas]. $D_{ES}=150$, $D_{EO}=120$, $D_{VS}=100$, $D_{MS}=110$.
\end{itemize}

\medskip
\textbf{Fracciones $\alpha_{gt}$:}\par\smallskip
\adjustbox{max width=\linewidth}{%
\renewcommand{\arraystretch}{1.3}%
\begin{tabular}{@{}lcc@{}}
\toprule
\textbf{Grupo $g$} & $\alpha_{g,d}$ (día) & $\alpha_{g,n}$ (noche)\\
\midrule
Esposa (ES)        & 0.30 & 0.30\\
Esposo (EO)        & 0.10 & 0.30\\
Varón soltero (VS) & 0.10 & 0.15\\
Mujer soltera (MS) & 0.10 & 0.20\\
\bottomrule
\end{tabular}}
\end{parambox}

\subsection{Variables de Decisión}

\begin{varbox}
\begin{itemize}[leftmargin=1.5em]
  \item $x_t \geq 0$: número de llamadas de tipo $t \in \mathcal{T}$ realizadas.
\end{itemize}
\end{varbox}

\subsection{Función Objetivo}

\begin{fobox}
Minimizar el costo total de la encuesta:
\[
  \min\; Z = \sum_{t\,\in\,\cset{T}} c_t\,x_t = 2x_d + 5x_n
\]
\end{fobox}

\subsection{Restricciones}

\begin{restbox}
\textbf{R1 — Cobertura mínima} de cada grupo $g \in \cset{G}$:
\[
  \sum_{t\,\in\,\cset{T}} \alpha_{gt}\,x_t \;\geq\; D_g
\]
Expandido:
\begin{align*}
  0.30x_d + 0.30x_n &\geq 150 \tag{R1-ES}\\
  0.10x_d + 0.30x_n &\geq 120 \tag{R1-EO}\\
  0.10x_d + 0.15x_n &\geq 100 \tag{R1-VS}\\
  0.10x_d + 0.20x_n &\geq 110 \tag{R1-MS}
\end{align*}

\textbf{R2 — Límite de llamadas nocturnas}
(a lo sumo la mitad del total):
\[
  x_n \leq \tfrac{1}{2}(x_d + x_n)
  \quad\Leftrightarrow\quad x_n \leq x_d
\]

\textbf{R3 — No negatividad}: $x_d, x_n \geq 0$.
\end{restbox}

\subsection{Modelo Matemático Completo}

\begin{modelbox}
\begin{alignat*}{2}
  \min\quad & Z = 2x_d + 5x_n\\[4pt]
  \text{s.a.}\quad
  & 0.30x_d + 0.30x_n \geq 150 &\quad& \text{(esposas)}\\
  & 0.10x_d + 0.30x_n \geq 120 &\quad& \text{(esposos)}\\
  & 0.10x_d + 0.15x_n \geq 100 &\quad& \text{(varones solteros)}\\
  & 0.10x_d + 0.20x_n \geq 110 &\quad& \text{(mujeres solteras)}\\
  & x_n \leq x_d               &\quad& \text{(límite nocturnas)}\\
  & x_d,\, x_n \geq 0          &\quad& \text{(no negatividad)}
\end{alignat*}
\end{modelbox}

\subsection{Solución Óptima}

Dado que $Z = 2400 - x_n$ sobre la recta de esposos (R1-EO activa), el costo
decrece con $x_n$ hasta que R1-MS impone $x_n \leq 100$:
\[
  \begin{cases} x_d + 3x_n = 1\,200 \\ x_d + 2x_n = 1\,100 \end{cases}
  \;\Rightarrow\; x_n^* = 100,\quad x_d^* = 900
\]

\begin{center}
\renewcommand{\arraystretch}{1.3}
\begin{tabular}{@{}lcccc@{}}
\toprule
\textbf{Restricción} & \textbf{LHS} & \textbf{RHS} & \textbf{Estado}\\
\midrule
R1-ES & 1\,000 & 500   & holgura\\
R1-EO & 1\,200 & 1\,200 & \textbf{activa}\\
R1-VS & 2\,100 & 2\,000 & holgura\\
R1-MS & 1\,100 & 1\,100 & \textbf{activa}\\
R2    & $100 \leq 900$ & — & holgura\\
\bottomrule
\end{tabular}
\end{center}

\begin{solucbox}
\[
  x_d^* = 900\text{ llamadas diurnas},\quad x_n^* = 100\text{ llamadas nocturnas}
\]
\[
  \boxed{Z^* = \$2\,300}
\]
\end{solucbox}

\begin{sensbox}
\begin{itemize}[leftmargin=1.5em]
  \item \textbf{R1-EO} (activa): precio sombra $= -\$1$/persona.
    Cada esposo adicional requerido cuesta \$1.
  \item \textbf{R1-MS} (activa): precio sombra $= -\$3$/persona.
    La restricción de mujeres solteras es 3 veces más costosa.
  \item \textbf{Costo nocturno $c_n$}: si $c_n$ baja a \$2, toda la encuesta podría
    realizarse de noche, reduciendo el costo significativamente.
  \item \textbf{Límite nocturnas R2}: holgura de 800 llamadas; no restrictiva.
\end{itemize}
\end{sensbox}

\subsection{Código AMPL}
\begin{lstlisting}[caption={Encuesta Telefónica -- Modelo AMPL}]
# -- Conjuntos -----------------------------------------
set G;   # grupos objetivo: {ES, EO, VS, MS}
set TL;  # tipos de llamada: {d, n}

# -- Parametros ----------------------------------------
param c     {TL};       # costo por llamada tipo t  [USD]
param alpha {G, TL};    # fraccion del grupo g alcanzada por llamada tipo t
param D     {G};        # detecciones minimas requeridas del grupo g

# -- Variables -----------------------------------------
var x {TL} >= 0;        # numero de llamadas tipo t

# -- Funcion objetivo ----------------------------------
minimize Costo:
    sum {t in TL} c[t] * x[t];

# -- Restricciones -------------------------------------
subject to Cobertura {g in G}:     # R1
    sum {t in TL} alpha[g,t] * x[t] >= D[g];

subject to LimNocturnas:           # R2
    x['n'] <= x['d'];

# -- Datos ---------------------------------------------
data;

set G  := ES EO VS MS;
set TL := d n;

param c :=  d 2   n 5;

param alpha :    d     n  :=
           ES  0.30  0.30
           EO  0.10  0.30
           VS  0.10  0.15
           MS  0.10  0.20 ;

param D := ES 150  EO 120  VS 100  MS 110;

# -- Resolver ------------------------------------------
option solver cplex;
solve;

display x;
display Costo;
# x[d]=900  x[n]=100  Costo=2300 USD
\end{lstlisting}

\clearpage
%%─────────────────────────────────────────────────────────────────────────────
\section{Ejercicio 6 — Brady Corporation}
%%─────────────────────────────────────────────────────────────────────────────

\subsection{Conjuntos}

\begin{conjbox}
\begin{itemize}[leftmargin=1.5em]
  \item $\mathcal{F}$: fuentes de madera $\{1, 2, T\}$ (grado 1, grado 2, troncos propios).
\end{itemize}
\end{conjbox}

\subsection{Parámetros}

\begin{parambox}
\begin{itemize}[leftmargin=1.5em]
  \item $c_f$: costo total por pie$^3$ de entrada de fuente $f \in \mathcal{F}$ [\$/pie$^3$]. $c_1=7$, $c_2=11$, $c_T=9.5$.
  \item $r_f$: rendimiento de madera útil de fuente $f \in \mathcal{F}$ [pie$^3$/pie$^3$]. $r_1=0.70$, $r_2=0.90$, $r_T=0.80$.
  \item $\tau_f$: tiempo de secado por pie$^3$ de fuente $f \in \mathcal{F}$ [s/pie$^3$]. $\tau_1=1.2$, $\tau_2=0.8$, $\tau_T=1.3$.
  \item $\sigma_f$: capacidad de oferta de fuente $f \in \mathcal{F}$ [pie$^3$/sem.]. $\sigma_1=40\,000$, $\sigma_2=60\,000$, $\sigma_T=35\,000$.
  \item $D$: demanda semanal de tablones útiles [pie$^3$/sem.]. $D = 90\,000$.
  \item $T_{\max}$: tiempo de secado disponible [s/sem.]. $T_{\max} = 144\,000$.
\end{itemize}
\textit{Nota: $c_T = 2.5$ (flete) $+ 3$ (corte) $+ 4$ (secado) $= 9.5$ \$/pie$^3$}
\end{parambox}

\subsection{Variables de Decisión}

\begin{varbox}
\begin{itemize}[leftmargin=1.5em]
  \item $x_f \geq 0$: pies$^3$ de madera adquiridos / enviados al aserradero desde la fuente $f \in \mathcal{F}$.
\end{itemize}
\end{varbox}

\subsection{Función Objetivo}

\begin{fobox}
Minimizar el costo semanal total de abastecimiento:
\[
  \min\; Z = \sum_{f\,\in\,\cset{F}} c_f\,x_f = 7x_1 + 11x_2 + 9.5x_T
\]
\end{fobox}

\subsection{Restricciones}

\begin{restbox}
\textbf{R1 — Satisfacción de demanda}
(rendimiento útil de cada fuente debe cubrir los 90\,000 pie$^3$ requeridos):
\[
  \sum_{f\,\in\,\cset{F}} r_f\,x_f \;\geq\; D
  \quad\Leftrightarrow\quad
  0.7x_1 + 0.9x_2 + 0.8x_T \geq 90\,000
\]

\textbf{R2 — Capacidad de oferta} por fuente $f$:
\[
  x_f \leq \sigma_f, \quad f \in \cset{F}
  \quad\Rightarrow\quad
  x_1 \leq 40\,000,\;\; x_2 \leq 60\,000,\;\; x_T \leq 35\,000
\]

\textbf{R3 — Tiempo de secado disponible}:
\[
  \sum_{f\,\in\,\cset{F}} \tau_f\,x_f \;\leq\; T_{\max}
  \quad\Leftrightarrow\quad
  1.2x_1 + 0.8x_2 + 1.3x_T \leq 144\,000
\]

\textbf{R4 — No negatividad}: $x_f \geq 0 \quad \forall\,f \in \cset{F}$.
\end{restbox}

\subsection{Modelo Matemático Completo}

\begin{modelbox}
\begin{alignat*}{2}
  \min\quad & Z = 7x_1 + 11x_2 + 9.5x_T\\[4pt]
  \text{s.a.}\quad
  & 0.7x_1 + 0.9x_2 + 0.8x_T \geq 90\,000 &\quad& \text{(demanda)}\\
  & x_1 \leq 40\,000 &\quad& \text{(oferta grado 1)}\\
  & x_2 \leq 60\,000 &\quad& \text{(oferta grado 2)}\\
  & x_T \leq 35\,000 &\quad& \text{(capacidad aserradero)}\\
  & 1.2x_1 + 0.8x_2 + 1.3x_T \leq 144\,000 &\quad& \text{(tiempo secado)}\\
  & x_1, x_2, x_T \geq 0 &\quad& \text{(no negatividad)}
\end{alignat*}
\end{modelbox}

\subsection{Solución Óptima}

Costo por pie$^3$ útil: $c_1/r_1 = 10 < c_T/r_T = 11.88 < c_2/r_2 = 12.22$.
Se maximiza la fuente más barata primero:

\begin{center}
\renewcommand{\arraystretch}{1.3}
\begin{tabular}{@{}lcccc@{}}
\toprule
\textbf{Fuente $f$} & $x_f^*$ & \textbf{Útil} (pie$^3$) & \textbf{Costo} & \textbf{Tiempo} (s)\\
\midrule
Grado 1 & 40\,000 & 28\,000 & \$280\,000 & 48\,000\\
Troncos & 35\,000 & 28\,000 & \$332\,500 & 45\,500\\
Grado 2 & 37\,778 & 34\,000 & \$415\,556 & 30\,222\\
\midrule
\textbf{Total} & & \textbf{90\,000} & \textbf{\$1\,028\,056} & 123\,722 $\leq$ 144\,000 \checkmark\\
\bottomrule
\end{tabular}
\end{center}

\begin{solucbox}
\[
  x_1^* = 40\,000,\quad x_2^* = 37\,778,\quad x_T^* = 35\,000\;\text{pies}^3
\]
\[
  \boxed{Z^* = \$1\,028\,056/\text{semana}}
\]
Restricciones activas: R1 (demanda), R2-G1 (oferta grado 1) y R2-T (aserradero).
\end{solucbox}

\begin{sensbox}
\begin{itemize}[leftmargin=1.5em]
  \item \textbf{Oferta grado 1 R2-G1} (activa): precio sombra $\approx -\$2.86$/pie$^3$.
    Cada pie$^3$ adicional de grado 1 disponible reduce el costo en \$2.86.
  \item \textbf{Aserradero R2-T} (activo): precio sombra $\approx -\$0.86$/pie$^3$.
    Ampliar la capacidad del aserradero en 1\,000 pie$^3$ ahorraría $\approx \$860$.
  \item \textbf{Tiempo de secado R3}: holgura $= 144\,000 - 123\,722 = 20\,278$ s ($\approx 5.6$ h); no restrictivo.
  \item \textbf{Oferta grado 2}: holgura $= 60\,000 - 37\,778 = 22\,222$ pie$^3$.
\end{itemize}
\end{sensbox}

\subsection{Código AMPL}
\begin{lstlisting}[caption={Brady Corporation -- Modelo AMPL}]
# -- Conjuntos -----------------------------------------
set F;              # fuentes: {g1, g2, troncos}

# -- Parametros ----------------------------------------
param c     {F};    # costo total por pie^3 de entrada   [USD/pie^3]
param r     {F};    # rendimiento de madera util          [pie^3/pie^3]
param tau   {F};    # tiempo de secado por pie^3          [s/pie^3]
param sigma {F};    # capacidad maxima de oferta          [pie^3/sem]
param D;            # demanda semanal de tablones utiles  [pie^3/sem]
param Tmax;         # tiempo total de secado disponible   [s/sem]

# -- Variables -----------------------------------------
var x {F} >= 0;     # pie^3 adquiridos/enviados de fuente f

# -- Funcion objetivo ----------------------------------
minimize Costo:
    sum {f in F} c[f] * x[f];

# -- Restricciones -------------------------------------
subject to Demanda:              # R1
    sum {f in F} r[f] * x[f] >= D;

subject to Oferta {f in F}:      # R2
    x[f] <= sigma[f];

subject to TiempoSecado:         # R3
    sum {f in F} tau[f] * x[f] <= Tmax;

# -- Datos ---------------------------------------------
data;

set F := g1 g2 troncos;

param c     := g1  7      g2 11      troncos  9.5;
param r     := g1  0.7    g2  0.9    troncos  0.8;
param tau   := g1  1.2    g2  0.8    troncos  1.3;
param sigma := g1 40000   g2 60000   troncos 35000;
param D    := 90000;
param Tmax := 144000;    # 40 h * 3600 s/h

# -- Resolver ------------------------------------------
option solver cplex;
solve;

display x;
display Costo;
# x[g1]=40000  x[g2]=37778  x[troncos]=35000
# Costo = 1028056 USD/semana
\end{lstlisting}

\clearpage
%%─────────────────────────────────────────────────────────────────────────────
\section{Ejercicio 7 — Centro de Reciclaje}
%%─────────────────────────────────────────────────────────────────────────────

\subsection{Conjuntos}

\begin{conjbox}
\begin{itemize}[leftmargin=1.5em]
  \item $\mathcal{M}$: tipos de chatarra $\{A, B\}$.
  \item $\mathcal{C}$: componentes de la aleación $\{\text{Al}, \text{Si}, \text{Ca}\}$ (aluminio, silicio, carbón).
\end{itemize}
\end{conjbox}

\subsection{Parámetros}

\begin{parambox}
\begin{itemize}[leftmargin=1.5em]
  \item $c_m$: costo por tonelada de chatarra $m \in \mathcal{M}$ [\$/t]. $c_A=100$, $c_B=80$.
  \item $a_{km}$: fracción del componente $k \in \mathcal{C}$ en la chatarra $m \in \mathcal{M}$.
  \item $\underline{q}_k$: contenido mínimo del componente $k \in \mathcal{C}$ en la aleación [t].
  \item $\overline{q}_k$: contenido máximo del componente $k \in \mathcal{C}$ en la aleación [t].
  \item $Q$: producción total de aleación requerida [t]. $Q = 1\,000$.
\end{itemize}

\medskip
\begin{center}
\renewcommand{\arraystretch}{1.3}
\begin{tabular}{@{}lcccc@{}}
\toprule
\textbf{Componente $k$} & $a_{k,A}$ & $a_{k,B}$ & $\underline{q}_k$ (t) & $\overline{q}_k$ (t)\\
\midrule
Aluminio (Al) & 0.06 & 0.03 & 30 & 60\\
Silicio  (Si) & 0.03 & 0.06 & 30 & 50\\
Carbón   (Ca) & 0.04 & 0.03 & 30 & 70\\
\bottomrule
\end{tabular}
\end{center}
\end{parambox}

\subsection{Variables de Decisión}

\begin{varbox}
\begin{itemize}[leftmargin=1.5em]
  \item $x_m \geq 0$: toneladas de chatarra $m \in \mathcal{M}$ utilizadas en la mezcla.
\end{itemize}
\end{varbox}

\subsection{Función Objetivo}

\begin{fobox}
Minimizar el costo total de las chatarras para producir 1\,000 toneladas de aleación:
\[
  \min\; Z = \sum_{m\,\in\,\cset{M}} c_m\,x_m = 100x_A + 80x_B
\]
\end{fobox}

\subsection{Restricciones}

\begin{restbox}
\textbf{R1 — Producción total exacta}:
\[
  \sum_{m\,\in\,\cset{M}} x_m = Q
  \quad\Leftrightarrow\quad x_A + x_B = 1\,000
\]

\textbf{R2 — Especificaciones de composición}
(para cada componente $k \in \cset{C}$, el contenido total debe
estar entre $\underline{q}_k$ y $\overline{q}_k$ toneladas):
\[
  \underline{q}_k \;\leq\; \sum_{m\,\in\,\cset{M}} a_{km}\,x_m \;\leq\; \overline{q}_k
\]
Expandido:
\begin{alignat*}{3}
  30 &\leq 0.06x_A + 0.03x_B &\leq 60 \quad& \text{(aluminio)}\\
  30 &\leq 0.03x_A + 0.06x_B &\leq 50 \quad& \text{(silicio)}\\
  30 &\leq 0.04x_A + 0.03x_B &\leq 70 \quad& \text{(carbón)}
\end{alignat*}

\textbf{R3 — No negatividad}: $x_A, x_B \geq 0$.
\end{restbox}

\subsection{Modelo Matemático Completo}

\begin{modelbox}
\begin{alignat*}{2}
  \min\quad & Z = 100x_A + 80x_B\\[4pt]
  \text{s.a.}\quad
  & x_A + x_B = 1\,000 &\quad& \text{(producción total)}\\
  & 0.06x_A + 0.03x_B \geq 30 &\quad& \text{(Al mín 3\%)}\\
  & 0.06x_A + 0.03x_B \leq 60 &\quad& \text{(Al máx 6\%)}\\
  & 0.03x_A + 0.06x_B \geq 30 &\quad& \text{(Si mín 3\%)}\\
  & 0.03x_A + 0.06x_B \leq 50 &\quad& \text{(Si máx 5\%)} \leftarrow\text{activa}\\
  & 0.04x_A + 0.03x_B \geq 30 &\quad& \text{(Ca mín 3\%)}\\
  & 0.04x_A + 0.03x_B \leq 70 &\quad& \text{(Ca máx 7\%)}\\
  & x_A,\, x_B \geq 0         &\quad& \text{(no negatividad)}
\end{alignat*}
\end{modelbox}

\subsection{Solución Óptima}

Sustituyendo $x_B = 1\,000 - x_A$ en cada restricción, la única que acota $x_A$ por abajo
es la de silicio máximo:
\[
  0.03x_A + 0.06(1\,000-x_A) \leq 50
  \;\;\Rightarrow\;\;
  60 - 0.03x_A \leq 50
  \;\;\Rightarrow\;\;
  x_A \geq \frac{1\,000}{3}
\]

Como $Z = 20x_A + 80\,000$ es creciente en $x_A$, se toma el mínimo factible:

\begin{center}
\renewcommand{\arraystretch}{1.3}
\begin{tabular}{@{}lcccc@{}}
\toprule
\textbf{Comp.} & \textbf{Contenido (t)} & \textbf{\%} & \textbf{Rango (\%)} & \textbf{Estado}\\
\midrule
Aluminio & 40.00 & 4.0 & [3, 6] & holgura\\
Silicio  & 50.00 & 5.0 & [3, 5] & \textbf{activa (máx)}\\
Carbón   & 33.33 & 3.3 & [3, 7] & holgura\\
\bottomrule
\end{tabular}
\end{center}

\begin{solucbox}
\[
  x_A^* = \frac{1\,000}{3} \approx 333.33\;\text{t},\quad
  x_B^* = \frac{2\,000}{3} \approx 666.67\;\text{t}
\]
\[
  \boxed{Z^* = \$86\,666.67}
\]
\end{solucbox}

\begin{sensbox}
\begin{itemize}[leftmargin=1.5em]
  \item \textbf{R2-Si máx} (activa): precio sombra $= -\$666.67$/t.
    Relajar el límite de silicio a 5.1\% (10 t más) reduce el costo $\approx\$667$.
  \item \textbf{Costo chatarra A}: la solución es óptima mientras $c_A > c_B = \$80$.
    Si $c_A \leq \$80$, se usaría solo chatarra A ($x_A = 1\,000$ t).
  \item \textbf{Rango de optimalidad}: $x_A^* = 1000/3$ es óptimo mientras
    el ratio $c_A/c_B \in (1, +\infty)$ con la restricción de silicio activa.
  \item Todas las demás restricciones tienen holgura amplia.
\end{itemize}
\end{sensbox}

\subsection{Código AMPL}
\begin{lstlisting}[caption={Centro de Reciclaje -- Modelo AMPL}]
# -- Conjuntos -----------------------------------------
set M;          # chatarras:   {A, B}
set C;          # componentes: {Al, Si, Ca}

# -- Parametros ----------------------------------------
param costo {M};        # costo por tonelada de chatarra m  [USD/t]
param a     {C, M};     # fraccion del componente k en chatarra m
param qmin  {C};        # contenido minimo del componente k  [t]
param qmax  {C};        # contenido maximo del componente k  [t]
param Q;                # produccion total requerida          [t]

# -- Variables -----------------------------------------
var x {M} >= 0;         # toneladas de chatarra m utilizadas

# -- Funcion objetivo ----------------------------------
minimize Costo:
    sum {m in M} costo[m] * x[m];

# -- Restricciones -------------------------------------
subject to Produccion:                  # R1
    sum {m in M} x[m] = Q;

subject to CompMin {k in C}:            # R2 (limite inferior)
    sum {m in M} a[k,m] * x[m] >= qmin[k];

subject to CompMax {k in C}:            # R2 (limite superior)
    sum {m in M} a[k,m] * x[m] <= qmax[k];

# -- Datos ---------------------------------------------
data;

set M := A B;
set C := Al Si Ca;

param costo := A 100   B 80;

param a :      A      B    :=
         Al  0.06   0.03
         Si  0.03   0.06
         Ca  0.04   0.03  ;

param qmin := Al 30   Si 30   Ca 30;
param qmax := Al 60   Si 50   Ca 70;
param Q    := 1000;

# -- Resolver ------------------------------------------
option solver cplex;
solve;

display x;
display Costo;
# x[A]=333.33 t   x[B]=666.67 t   Costo=86666.67 USD
\end{lstlisting}

\clearpage
%%─────────────────────────────────────────────────────────────────────────────
\section*{Resumen General}
\addcontentsline{toc}{section}{Resumen General}
%%─────────────────────────────────────────────────────────────────────────────

\begin{center}
\renewcommand{\arraystretch}{1.5}
\begin{tabular}{@{}clccc@{}}
\toprule
\textbf{\#} & \textbf{Ejercicio} & \textbf{Tipo} & \textbf{Restricciones activas} & \textbf{Óptimo}\\
\midrule
1 & La Química       & min & Demanda A, Demanda C       & \$450/día\\
2 & Charcha          & max & Insumos 1--2, Río, Tiempo  & \$6\,366/semana\\
3 & Refinería Melrose & max & Cap.\ gas/oil, Calidad     & \$11\,230\\
4 & Donovan          & min & Cap.\ T2, T3, T4           & \$495\,000/año\\
5 & Encuesta         & min & Esposos, Muj.\,solteras    & \$2\,300\\
6 & Brady            & min & Oferta G1, Aserradero      & \$1\,028\,056/semana\\
7 & Reciclaje        & min & Silicio máx 5\%            & \$86\,667/lote\\
\bottomrule
\end{tabular}
\end{center}

\end{document}
